% !TeX root = ../beamer.tex

\section{Text and boxes}

\subsection{Outline slides}
% Steps on the left, the outline they produce on the right
\begin{frame}{Outline slides}
\begin{columns}[T,onlytextwidth]
\begin{column}{0.5\textwidth}
\step{1. Choose when it appears}
\begin{codeonly}
\usetheme[sectionslides]
   {SimplePlusNord}
\end{codeonly}
\smallskip
\step{2. List the slides}
\begin{codeonly}
\subsection{Title}
\end{codeonly}
\smallskip
\step{3. Optional}
\begin{codeonly}
\tocnumbers
\end{codeonly}
\end{column}
\begin{column}{0.45\textwidth}
% no \stepoffset, and the extra section is left out: the list is long and would not fit
\tableofcontents[currentsection,sections={1-3},subsectionstyle=show/show/hide]
\end{column}
\end{columns}
\end{frame}

\subsection*{Outside the outline}
% this slide is not in the outline: its subsection is starred
\begin{frame}{Outside the outline}
\begin{columns}[T,onlytextwidth]
\begin{column}{0.5\textwidth}
\step{1. Star the subsection}
\begin{codeonly}
\subsection*{Backup}
\end{codeonly}
\smallskip
\step{2. Or write none}
\begin{codeonly}
% no \subsection before
% the slide: it joins the
% entry that comes before
\end{codeonly}
\end{column}
\begin{column}{0.45\textwidth}
\stepoffset
This slide is the first case: it is not in the outline.

\medskip
The same star on a section hides the whole section.
\end{column}
\end{columns}
\end{frame}

\subsection{Extra section}
\begin{frame}{Extra section}
\begin{columns}[T,onlytextwidth]
\begin{column}{0.5\textwidth}
\step{1. Use it for the references}
\begin{codeonly}
\extrasection{References}
\end{codeonly}
\end{column}
\begin{column}{0.45\textwidth}
\stepoffset
It takes the place of a normal section. The outline keeps it apart: last, in blue, as a line under the columns, and with no slide of its own to open it.

\medskip
Look at the outline slides: it is "References".
\end{column}
\end{columns}
\end{frame}

\subsection{Bullet Points}
\begin{frame}{Bullet Points}
    \begin{itemize}
        \item A short point
        \item A second point, a little longer than the first
        \item A third point, long enough to wrap onto a second line, so that the hanging indent of the list shows
        \item A fourth point
        \item A fifth point, with nothing more to say
        \item \fade{Faded with \texttt{\textbackslash fade}, \textcolor{nord10}{and the blue inside is lightened too}}
    \end{itemize}
\end{frame}

\subsection{Step by step}
% [c] changes nothing: it keeps the way of the frame options under test
\begin{frame}[c]{Step by step}
\begin{columns}[T,onlytextwidth]
\begin{column}{0.5\textwidth}
\step{1. Reveal with a pause}
\begin{codeonly}
First point.
\pause
Second point.
\source{Verdi (2021)}
\end{codeonly}
\smallskip
\step{2. Code in a step}
\begin{codeonly}
\begin{onlyenv}<2>
  \begin{codeblock}{Python}
  ...
  \end{codeblock}
\end{onlyenv}
\end{codeonly}
\end{column}
\begin{column}{0.45\textwidth}
\stepoffset
First point.
\pause

Second point, with its source.
\source{Verdi (2021)}
\begin{onlyenv}<2>
\begin{codeblock}{Python}
print("second step")
\end{codeblock}
\end{onlyenv}
\end{column}
\end{columns}
\end{frame}

\subsection{Multiple Columns}
\begin{frame}{Multiple Columns}
    \begin{columns}[c]
        \column{.45\textwidth}
        \textbf{Steps}
        \begin{enumerate}
            \item State the claim
            \item Explain it
            \item Give an example
        \end{enumerate}

        \column{.45\textwidth}
        Columns split the slide in parts of the width you choose, and each part flows as its own text, as in a newspaper. They are centred vertically here; \texttt{[T]} aligns them on their first line.
    \end{columns}
\end{frame}

\subsection{Boxes of Highlighted Text}
\begin{frame}{Boxes of Highlighted Text}
    A few words can stand out in \alert{bold}; use it sparingly.

    \begin{infobox}{Box}
        Sample text
    \end{infobox}

    \begin{alertbox}{Alertbox}
        Sample text in a box that warns
    \end{alertbox}

    \begin{examplebox}{Examplebox}
        Sample text in a box that illustrates
    \end{examplebox}
\end{frame}

\subsection{Boxes: code and result}
% Steps with their code on the left, the result on the right
\begin{frame}{Boxes: code and result}
\begin{columns}[T,onlytextwidth]
\begin{column}{0.5\textwidth}
\step{1. A box}
\begin{codeonly}
\begin{infobox}{Title}
  Text
\end{infobox}
\end{codeonly}
\end{column}
\begin{column}{0.45\textwidth}
\stepoffset
\begin{infobox}{Box}
    Sample text
\end{infobox}
\end{column}
\end{columns}
\bigskip
\begin{columns}[T,onlytextwidth]
\begin{column}{0.5\textwidth}
\step{2. Alert and example}
\begin{codeonly}
\begin{alertbox}{Title}
\begin{examplebox}{Title}
\end{codeonly}
\end{column}
\begin{column}{0.45\textwidth}
\stepoffset
\begin{alertbox}{Alertbox}
    Sample text
\end{alertbox}
\begin{examplebox}{Examplebox}
    Sample text
\end{examplebox}
\end{column}
\end{columns}
\end{frame}
